\chapter{研究方法与实验设计}
\label{cha:thirdsection}

\section{实验一：训练样本数量影响}
实验选取训练点数分别为15,35,55,75,100


\section{实验二：网络结构影响}

\section{实验三：激活函数影响}
\begin{table}[htbp]
  \centering
  \caption{模型默认参数}
    \begin{tabular}{ccccc}
    \toprule
    激活函数 & MSE & MAE \\
    \midrule
        relu & 0.002412 & 0.039204 \\
        tanh & 0.000025 & 0.002281 \\
        sigmoid & 0.000091 & 0.005753 \\ 
    \bottomrule
    \end{tabular}%
  \label{Tab4-1}%
\end{table}%

结果表明，随着训练样本数量增加，预测误差逐渐减小，模型的拟合能力明显增强；而Loss曲线显示，样本数量越多，收敛越稳定